\section{Relationship Between Cost and Error}
\label{sec:cost_and_error}

In \autoref{sec:setup}, we claim that there is a tight connection
between reducing any cost function $\fairconst{h_t}{t}$ and reducing the generalized
error rates $\fp{h_t}$ and $\fn{h_t}$
for approximately calibrated classifiers.
In other words, assuming we are approximately calibrated, improving cost will
approximately improve our error rates.
We formalize this notion in this section:
%
\begin{lemma}
  \label{lma:cost_error}
  Let $h_t$ be a classifier with $\epsilon(h_t) = \delta_{cal}$ and cost
  $\fairconst{h_t}{t}$. For any other classifier $h_t'$, if $\fp{h_t'} <
  \fp{h_t} - \frac{4\delta_{cal}}{1-\mu_t}$ or $\fn{h_t'} < \fn{h_t} -
  \frac{4\delta_{cal}}{\mu_t}$, then $\fairconst{h_t'}{t} <
  \fairconst{h_t}{t}$ or $\epsilon(h_t) > \delta_{cal}$.
  \label{lem:fpfn}
\end{lemma}
\begin{proof}
  First, assume that $\fp{h_t'} < \fp{h_t} - \frac{4\delta_{cal}}{1-\mu_t}$.
  Then, there are two cases: either $\fn{h_t'} < \fn{h_t}$ or $\fn{h_t'} \ge
  \fn{h_t}$. In the first case, $\fairconst{h_t'}{t} < \fairconst{h_t}{t}$
  because $\fp{h_t'} < \fp{h_t}$ and $\fn{h_t'} < \fn{h_t}$. In the second case,
  if $\epsilon(h_t') \le \delta_{cal}$, we can use Lemma
  \ref{lem:approx-linear-cal} to get
  \begin{align*}
    \fp{h_t'}
    &\ge \frac{\mu_t}{1-\mu_t} \fn{h_t'} -
    \frac{2\delta_{cal}}{1-\mu_t} \\
    &\ge \frac{\mu_t}{1-\mu_t} \fn{h_t} -
    \frac{2\delta_{cal}}{1-\mu_t} \\
    &\ge \fp{h_t} -
    \frac{4\delta_{cal}}{1-\mu_t}
  \end{align*}
  Since this contradicts the initial assumption that $\fp{h_t'} < \fp{h_t} -
  \frac{4\delta_{cal}}{1-\mu_t}$, it cannot be the case that $\epsilon(h_t') \le
  \delta_{cal}$. This proves the lemma when $\fp{h_t'} < \fp{h_t} -
  \frac{4\delta_{cal}}{1-\mu_t}$.

  To prove the second part, we now assume that $\fn{h_t'} < \fn{h_t} -
  \frac{4\delta_{cal}}{\mu_t}$. We again break this into two cases. If
  $\fp{h_t'} < \fp{h_t}$, then $\fairconst{h_t'}{t} < \fairconst{h_t}{t}$. If
  $\fp{h_t'} \ge \fp{h_t}$, then under the assumption that $\epsilon(h_t') \le
  \delta_{cal}$, we can rearrange Lemma \ref{lem:approx-linear-cal} to get
  \begin{align*}
    \fn{h_t'}
    &\ge \frac{1-\mu_t}{\mu_t} \fp{h_t'} -
    \frac{2\delta_{cal}}{\mu_t} \\
    &\ge \frac{1-\mu_t}{\mu_t} \fp{h_t} -
    \frac{2\delta_{cal}}{\mu_t} \\
    &\ge \fn{h_t} -
    \frac{4\delta_{cal}}{\mu_t}
  \end{align*}
  Again, this contradicts the assumption that $\fn{h_t'} < \fn{h_t} -
  \frac{4\delta_{cal}}{\mu_t}$, so it cannot be the case that
  $\epsilon(h_t') \le \delta_{cal}$. This completes the proof.
\end{proof}
%
From this result we can derive a stronger claim
for perfectly calibrated classifiers.
%
\begin{lemma}
  Let $h_t$ and $h'_t$ be perfectly calibrated classifiers with cost
  $\fairconst{h_t}{t} \leq \fairconst{h'_t}{t}$ for some
  cost function $\fairconstname$. Then $\fp{h_t} \leq \fp{h'_t}$
  and $\fn{h_t} \leq \fn{h'_t}$, with equality only if $\fairconst{h_t}{t} = \fairconst{h'_t}{t}$.
\end{lemma}
